\section{Introduction}

Studying the similarity among closely related genomes is a major task
in comparative genomics. Conventionally, genomes are globally aligned
to identify common homologous regions, which can then be analyzed
further, for example, to calculate phylogenetic distances or perform
functional annotation. However, aligning multiple genomes can be both
resource-intensive and time-consuming, especially when dealing with
entire genomes.

A global alignment between two sequences can be approximated as a
series of long, gapless exact matches (called anchors) separated by
mismatches. Using enhanced suffix arrays (ESA) to find these exact
matches makes such an approximation fast compared to an actual
alignment (Figure~\ref{fig:Intro1}).

\begin{figure}[H]
  \centering
    \includegraphics[width=0.5\linewidth]{figIntro1.eps}
    \caption{Conceptual comparison of a global alignment (GA) of two
      sequences ($S_1$ and $S_2$) and anchor-based alignment
      approximation (AA). The GA of these sequences can be understood
      as a process that aims to maximize the coverage of $S_2$ with
      nucleotides from the sequence $S_1$ that match. GA contains gaps
      shown with a thinner line. In contrast, AA approximates the
      alignment using a series of gapless 'local alignments' that
      consist of exact matches shown as regions of $S_2$ separeted
      with mismatches (mutations) shown in darker grey. While AA is
      not identical to GA, it does capture the key similarities
      (\textit{homologous regions}), between the sequences.}
    \label{fig:Intro1}
\end{figure}

Such an anchor-based approach has been implemented in various
programs, such as \ty{andi}~\cite{hau15:and} and
\ty{phylonium}~\cite{kloe19:phy}, which are used for the quick
calculation of evolutionary distances. The program
\texttt{fur}~\cite{hau21:fur} also uses this method in its
Intersection step to identify common regions across multiple genomes.

The \ty{chr} package, which stands for 'common homologous regions' is
designed to approximate global multi-sequence alignments using
exact-match-based anchors. Its primary function is \ty{Intersect},
which is used to find common homologous regions in multiple similar
genomes.
\newpage
